\section{Introduction}
\label{sec:introduction}

Logarithmic averaging is often substantially more accessible than
ordinary averaging in multiplicative number theory.  In particular,
Tao proved logarithmically averaged two-point Chowla and Elliott
theorems for fixed affine forms \citep{Tao2016LogChowla}.  The
literal fixed-shift route considered in TPC-32 and TPC-33, however,
returns ordinary sums over a growing family of affine columns, with a
prescribed nonzero determinant \(h_0\), actual masks, short fibers,
and a complete outer reassembly \citep{WangTPC32,WangTPC33}.

TPC-85 exhibited a bounded dyadic sequence showing that no
structure-free black box simply deletes the logarithmic weight
\citep{WangTPC85}.  The purpose of the present paper is to sharpen
that obstruction into an exact theorem:
\begin{quote}
the missing datum is the signed Ces\`aro defect of the
\emph{unnormalized} logarithmic primitive.
\end{quote}
This distinction matters.  If
\[
 H_N=\sum_{n\le N}\frac1n,\qquad
 \lambda_N=\frac1{H_N}\sum_{n\le N}\frac{a_n}{n},
\]
then \(\lambda_N\) is the normalized logarithmic mean, while
\[
 L_N=H_N\lambda_N
\]
is the unnormalized logarithmic primitive.  The exact discrete Abel
identity is written in terms of \(L_N\), not \(\lambda_N\).

TPC-86 develops a different positive interface: coherent Fourier
reassembly of the distinguished zero mode from a controlled
low-frequency window \citep{WangTPC86}.  The present paper does not
replace that route.  It audits the separate possibility of obtaining
ordinary affine cancellation from logarithmically averaged
correlation information.

\subsection{Main results}

For an arbitrary complex sequence \(a=(a_n)\), define
\[
 S_N=\sum_{n=1}^Na_n,\qquad
 L_N=\sum_{n=1}^N\frac{a_n}{n},\qquad S_0=L_0=0.
\]
Our first result is the pointwise identity
\begin{equation}
 \frac{S_N}{N}
 =L_N-\frac1N\sum_{m=0}^{N-1}L_m
 =\frac1N\sum_{m=0}^{N-1}(L_N-L_m).
\label{eq:intro-abel}
\end{equation}
Therefore \(S_N/N\to0\) if and only if the signed quantity on the
right tends to zero.  This is not merely an implication: the two
objects are equal for every \(N\), so there is no hidden exponent
loss.

Convergence \(L_N\to\ell\) is sufficient.  Boundedness \(L_N=O(1)\)
is sharply insufficient.  The dyadic block sequence
\[
 a_n=(-1)^{\lfloor\log_2n\rfloor}
\]
has a bounded logarithmic primitive and normalized logarithmic mean
\(\lambda_N\to0\), yet its natural means oscillate between
\(\pm1/3\) along complete dyadic blocks.  More strongly, every
terminal multiplicative interval has uniformly bounded logarithmic
sum, so every Tao-shaped normalized window with
\(\omega(x)\to\infty\) has cancellation.

The exact criterion also identifies useful additional hypotheses.
We prove:
\begin{enumerate}[label=\textup{(\roman*)},leftmargin=2.5em]
 \item an absolute Ces\`aro oscillation criterion;
 \item a terminal-half local criterion
 \[
  \max_{N/2\le m\le N}
  \left|\sum_{m<n\le N}\frac{a_n}{n}\right|\to0;
 \]
 \item a rate-preserving version for fixed powers below one;
 \item an explicit dyadic variation convolution bound.
\end{enumerate}
These are genuine Tauberian inputs: none follows from boundedness of
the global logarithmic primitive alone.

\subsection{Arithmetic status}

The results have three proof levels.
\begin{enumerate}[label=\textup{\(\mathrm{L}\arabic*\):},leftmargin=3.2em,start=0]
 \item finite discrete identities, quantitative inequalities, and
 structure-free countermodels;
 \item a literal transfer certificate retaining the same fixed
 \(h_0\), growing affine coefficients, actual fibers, masks, outer
 reassembly, and global normalization;
 \item a new uniform M\"obius-correlation estimate for that actual
 growing family.
\end{enumerate}
This paper proves \(\Lzero\) results and formulates the \(\Lone\)
certificate.  It proves no \(\Ltwo\) estimate.  In particular, the
dyadic sequence is not multiplicative and does not obstruct a proof
that uses M\"obius-specific structure.

\subsection{Organization}

\Cref{sec:discrete-identity} proves the exact Abel identities.
\Cref{sec:cesaro-criterion} gives the necessary-and-sufficient
criterion and the convergence theorem.  Quantitative local and
variation conditions appear in \cref{sec:local-variation}.
\Cref{sec:dyadic-counterexample} proves the sharp dyadic
counterexample.  The theorem-quantifier audit is
\cref{sec:literature-quantifiers}; the literal physical certificate
and endpoint ledger are in \cref{sec:physical-interface}.  The claim
boundary is recorded in \cref{sec:claim-boundary}.
